\section{Finitud y control uniforme de las longitudes cohomológicas}
\label{sec:finitud-uniforme}

La Nota~\ref{sec:construccion-selmer} construyó clases de Selmer y
calculó pérdidas locales. Para insertar sus longitudes en un presupuesto
de denominadores hace falta controlar los grupos completos. Esta nota
resuelve primero una cuestión cualitativa: para un cuerpo totalmente
imaginario fijo y un corte finito de torcidos, el presupuesto
$\mathcal V(E,t)$ de la Nota~\ref{sec:cohomologia} es finito. La entrada
clásica es la comparación entre cohomología y los grupos pares de
teoría $K$. Sus cardinalidades proporcionan, además, una expresión
exacta del presupuesto una vez descontada la corrección de dualidad.

Después calculamos esa corrección, las longitudes de conúcleos
presentados por matrices y el crecimiento por capas con torsión.
Una suma factorial mejora el coste de los torcidos alejados de los
centros de C. Estos resultados precisan qué se puede controlar
uniformemente y qué información aritmética sigue haciendo falta
cuando cambian los cuerpos y el corte.

\fuente{Nota~\ref{sec:cohomologia} y C, §§2--6, pp.~2--9,
para las dos cotas de torre; C, §§7--8, pp.~9--15, para su puente
aritmético. A93, Proposición~9.4 y desarrollo, pp.~67--74,
distingue la finitud de la cota cuantitativa \texttt{LemmaK}.
Weibel~\cite{weibel-kbook}, IV.1.17--1.18 y IV.6.9,
VI.8.2 y VI.8.4, aporta las entradas clásicas de teoría $K$.
Las deducciones de longitudes, invariantes, imágenes y sumas
se prueban aquí; no se añade una lectura ni una ejecución de Lean.}

\subsection{El presupuesto y la entrada clásica}

Sea $E$ un cuerpo de números totalmente imaginario. Para un primo
$\ell$ y un entero $e\ge1$, conservamos
\begin{equation}
 \begin{aligned}
 A_{\ell,e}&=(\Q_\ell/\Z_\ell)(-e),\\
 \operatorname{Sha}_{\ell,e}(E)
 &=\ker\!\left(H^1(G_{E,S_\ell},A_{\ell,e})
       \longrightarrow\prod_{v\mid\ell}H^1(E_v,A_{\ell,e})\right),\\
 s_{\ell,e}(E)&=\length_{\Z_\ell}\operatorname{Sha}_{\ell,e}(E),\qquad
 W_\ell(E,e)=\length_{\Z_\ell}H^2(G_{E,S_\ell},\Z_\ell(e+1)).
 \end{aligned}
 \label{eq:fu-grupos}
\end{equation}
$S_\ell$ consta de los lugares sobre $\ell$ y los infinitos;
los lugares complejos no añaden condiciones.
Todas las cohomologías son continuas.
Hasta justificar su finitud, las longitudes admiten el valor infinito.
Para $t\ge1$ entero, el presupuesto es
\begin{equation}
 \mathcal V(E,t)
 =\sup_{\mathcal P\text{ finito}}
   \sum_{\ell\in\mathcal P}\sum_{1\le e<t}
       s_{\ell,e}(E)\log\ell.
 \label{eq:fu-presupuesto}
\end{equation}
Esta definición no presupone convergencia de una serie.

\paragraph{La entrada de teoría $K$.}
Los grupos abelianos $K_n(\mathcal O_E)$ son finitamente generados
por el teorema de Quillen; el cálculo de rangos de Borel da rango
cero para $n=2e\ge2$. Por tanto $K_{2e}(\mathcal O_E)$ es finito.
La comparación cohomológica proporciona
\begin{equation}
 K_{2e}(\mathcal O_E)\{\ell\}
 \simeq H^2_{\mathrm{\acute et}}
       (\mathcal O_E[1/\ell],\Z_\ell(e+1))
 \simeq H^2(G_{E,S_\ell},\Z_\ell(e+1)).
 \label{eq:fu-k-comparacion}
\end{equation}
Aquí $B\{\ell\}$ significa la parte $\ell$-primaria de un grupo
finito $B$. Son resultados clásicos recibidos, no consecuencias
de un cálculo de matrices: véase Weibel, IV.1.17--1.18, p.~12,
IV.6.9, p.~59, y VI.8.2--8.4, pp.~47--48, con la paginación
de los PDF independientes de esos capítulos.

La última identificación usa Milne~\cite{milne-dualidad},
II, Proposición~2.9, pp.~170--172, para los coeficientes finitos.
Sus grupos $H^1$ son finitos por II, Teorema~2.13, p.~174;
el sistema inverso satisface la condición de Mittag--Leffler y
no introduce un término $\varprojlim{}^1 H^1$ al pasar a
$\Z_\ell(e+1)$. La hipótesis de cuerpo totalmente imaginario
permite incluir $\ell=2$ en \eqref{eq:fu-k-comparacion} sin
correcciones de lugares reales. No se exige que $E$ contenga
$\mu_\ell$ ni $\mu_4$, ni que sea abeliano.

\subsection{Finitud y soporte finito para un cuerpo fijo}

Definimos el número real finito
\begin{equation}
 \mathcal K(E,t)=\sum_{1\le e<t}\log\#K_{2e}(\mathcal O_E),
 \qquad \mathcal K(E,1)=0.
 \label{eq:fu-k-presupuesto}
\end{equation}
Es el logaritmo de un producto finito de enteros positivos.
Esta expresión no presupone que hayamos calculado esos enteros.

\begin{proposicion}[Presupuesto finito]
\label{prop:fu-finitud}
Para todo cuerpo totalmente imaginario $E$ y todo entero $t\ge1$,
los grupos $\operatorname{Sha}_{\ell,e}(E)$ y
$H^2(G_{E,S_\ell},\Z_\ell(e+1))$ son finitos y
\begin{equation}
 0\le\mathcal V(E,t)\le\mathcal K(E,t)<\infty.
 \label{eq:fu-v-menor-k}
\end{equation}
Para cada $1\le e<t$, si
$\operatorname{Sha}_{\ell,e}(E)\ne0$, entonces
$\ell\mid\#K_{2e}(\mathcal O_E)$. En particular, solo hay
finitos pares $(\ell,e)$ con $1\le e<t$ y longitud positiva.
\end{proposicion}

\begin{proof}
La entrada \eqref{eq:fu-k-comparacion} da la finitud de $H^2$ y
\[
 W_\ell(E,e)=v_\ell\bigl(\#K_{2e}(\mathcal O_E)\bigr).
\]
El fragmento de Poitou--Tate de la Nota~\ref{sec:cohomologia}
sobreyecta el dual de ese $H^2$ sobre
$\operatorname{Sha}_{\ell,e}(E)$. Por tanto este último es
finito y $0\le s_{\ell,e}(E)\le W_\ell(E,e)$.
Si $s_{\ell,e}(E)>0$, la valoración de la cardinalidad del
$K$-grupo es positiva. Para cada $e$ solo hay finitos primos
con esta propiedad. Sumando sobre cualquier conjunto finito
$\mathcal P$ y usando la factorización prima de cada cardinalidad,
\[
 \sum_{\ell\in\mathcal P}\sum_{1\le e<t}s_{\ell,e}(E)\log\ell
 \le\sum_{1\le e<t}\sum_\ell
     v_\ell\bigl(\#K_{2e}(\mathcal O_E)\bigr)\log\ell
 =\mathcal K(E,t).
\]
Tomar el supremo prueba la cota. Para $t=1$ ambas sumas son vacías.
\end{proof}

Se ha resuelto la sumabilidad para $E,t$ fijos mediante un soporte
aritmético finito. La afirmación es más fuerte que la finitud aislada
de cada grupo. Todavía no acota $\mathcal K(E,t)$ mediante el grado,
el discriminante o el corte; cambiar $E$ cambia los enteros que
soportan la suma.

\subsection{El cálculo de los invariantes locales}

Para una extensión finita $L/\Q_\ell$, escribimos
$\epsilon_L=e(L/\Q_\ell)$ para su índice de ramificación y ponemos
\[
 a_{\ell,e}(L)=\length_{\Z_\ell}H^0(L,A_{\ell,e}).
\]
La letra $e$ sigue designando el torcido. El número $\epsilon_L$
no es el grado residual.

\begin{lema}[Invariantes y ramificación]
\label{lem:fu-invariantes}
Los invariantes anteriores son finitos. Si $\ell$ es impar,
\begin{equation}
 \begin{gathered}
 a_{\ell,e}(L)>0\quad\Longrightarrow\quad
       \ell-1\mid e\epsilon_L,\\
 a_{\ell,e}(L)\le1+v_\ell(\epsilon_L)+v_\ell(e).
 \end{gathered}
 \label{eq:fu-h0-impar}
\end{equation}
En el primo diádico,
\begin{equation}
 a_{2,e}(L)\le2+v_2(\epsilon_L)+v_2(e).
 \label{eq:fu-h0-diadico}
\end{equation}
\end{lema}

\begin{proof}
Sea $H=\chi_\ell(G_L)\subseteq\Z_\ell^\times$ la imagen
ciclotómica. La extensión $\Q_\ell(\mu_{\ell^\infty})/\Q_\ell$
es totalmente ramificada. Su intersección con $L$ es finita y
\[
 d=[\Z_\ell^\times:H]
   =[L\cap\Q_\ell(\mu_{\ell^\infty}):\Q_\ell]
   \mid\epsilon_L.
\]
La divisibilidad usa la multiplicatividad de los índices de
ramificación en esa torre de cuerpos. En particular $H$ es abierto.

Un elemento de $\Q_\ell/\Z_\ell$ es invariante si lo anulan todos
los $u^{-e}-1$, con $u\in H$. Como
$u^{-e}-1=-u^{-e}(u^e-1)$, sus valoraciones son las mismas.
Si el ideal generado por los $u^e-1$ es $(\ell^b)$, el grupo
de invariantes es el subgrupo anulado por $\ell^b$ y tiene longitud
$b$. El ideal es no nulo porque una imagen abierta contiene
un elemento principal cuya potencia $e$ no es uno.

Para $\ell$ impar, la descomposición de las unidades y la estructura
de sus subgrupos cerrados dan
\[
 H=\mu_m\times(1+\ell^{b_0}\Z_\ell),\qquad
 m\mid\ell-1,\quad b_0\ge1,\quad
 d=\frac{\ell-1}{m}\ell^{b_0-1}.
\]
Si $m\nmid e$, algún elemento de $\mu_m$ tiene potencia $e$
distinta de uno módulo $\ell$; el ideal es la unidad y no hay
invariantes. Si $m\mid e$, el cálculo binomial
\[
 v_\ell\bigl((1+\ell^{b_0})^e-1\bigr)=b_0+v_\ell(e)
\]
y la continuidad sobre el grupo principal dan exactamente
$a_{\ell,e}(L)=b_0+v_\ell(e)$.
Ahora $b_0-1=v_\ell(d)\le v_\ell(\epsilon_L)$.
Además $(\ell-1)/m$ divide $\epsilon_L$, y $m\mid e$ implica
$\ell-1\mid e\epsilon_L$.

Para $\ell=2$ usamos $U=1+4\Z_2$. Su intersección con $H$ tiene
la forma $1+2^{b_0}\Z_2$, con $b_0\ge2$ y
\[
 2^{b_0-2}=[U:U\cap H]\mid d.
\]
El elemento $1+2^{b_0}$ basta para anular los invariantes después
de multiplicar por $2^{b_0+v_2(e)}$.
La identidad de valoración vale porque $b_0\ge2$.
Esto da $a_{2,e}(L)\le b_0+v_2(e)$ y la cota indicada.
La posible presencia de $-1$ puede reducir esa longitud;
no la hemos descartado para afirmar una igualdad.
\end{proof}

La identidad binomial utilizada se obtiene primero para un exponente
primo a $\ell$, y luego observando que elevar a $\ell$ aumenta la
valoración en uno. Para $\ell=2$ esa última afirmación exige nivel
al menos dos. Es la misma precaución que aparece en C, §4,
pp.~5--6, y A93, p.~71.

\subsection{La corrección de dualidad con una cota uniforme}

Escribimos también
$a_{\ell,e}(E)=\length H^0(G_{E,S_\ell},A_{\ell,e})$.
La restricción lo incluye en cada grupo local de invariantes;
por el lema es finito. El fragmento de dualidad utilizado en la
Nota~\ref{sec:cohomologia} es
\begin{equation}
 \begin{split}
 0\longrightarrow H^0(G_{E,S_\ell},A_{\ell,e})
 &\longrightarrow\prod_{v\mid\ell}H^0(E_v,A_{\ell,e})\\
 &\longrightarrow H^2(G_{E,S_\ell},\Z_\ell(e+1))^\vee
 \longrightarrow\operatorname{Sha}_{\ell,e}(E)
 \longrightarrow0.
 \end{split}
 \label{eq:fu-dualidad}
\end{equation}
La entrada de dualidad incluye el paso entre coeficientes finitos,
compactos y divisibles. El torcido dual de $e+1$ es $-e$.

\begin{proposicion}[Presupuesto exacto y error local]
\label{prop:fu-presupuesto-exacto}
Sea $E$ totalmente imaginario, $[E:\Q]\le D$, con $D\ge1$ entero,
y sea $t\ge2$. Pongamos $r=t-1$ y $X=Dr+1$. Entonces
\begin{equation}
 \begin{aligned}
 \mathcal K(E,t)-\mathcal V(E,t)
 &=\sum_\ell\sum_{e=1}^{r}(\log\ell)
       \left(\sum_{v\mid\ell}a_{\ell,e}(E_v)-a_{\ell,e}(E)\right),\\
 0\le\mathcal K(E,t)-\mathcal V(E,t)
 &\le DrX\log(2DrX).
 \end{aligned}
 \label{eq:fu-correccion-global}
\end{equation}
La suma de correcciones se soporta en los primos $\ell\le X$.
Para cada primo, con $c_\ell=1$ si $\ell$ es impar y $c_2=2$,
\begin{equation}
 0\le W_\ell(E,e)-s_{\ell,e}(E)
 \le D\bigl(c_\ell+\lfloor\log_\ell D\rfloor+v_\ell(e)\bigr).
 \label{eq:fu-correccion-primo}
\end{equation}
\end{proposicion}

\begin{proof}
Todos los términos de \eqref{eq:fu-dualidad} son finitos.
La aditividad de longitudes da
\begin{equation}
 W_\ell(E,e)-s_{\ell,e}(E)
 =\sum_{v\mid\ell}a_{\ell,e}(E_v)-a_{\ell,e}(E)\ge0.
 \label{eq:fu-defecto-local}
\end{equation}
Hay a lo sumo $D$ lugares sobre $\ell$ y cada índice de ramificación
es a lo sumo $D$. El lema anterior prueba
\eqref{eq:fu-correccion-primo} al omitir el término global,
que es no negativo.

Si $\ell$ es impar y algún invariante local es no nulo, entonces
$\ell-1\mid e\epsilon_{E_v}$, luego
$\ell\le e\epsilon_{E_v}+1\le Dr+1=X$.
El primo $2$ también es a lo sumo $X$. Para los términos no nulos,
\[
 \ell^{a_{\ell,e}(E_v)}\le
 \begin{cases}
  \ell\epsilon_{E_v}e\le DrX,&\ell\ne2,\\
  4\epsilon_{E_v}e\le4Dr\le2DrX,&\ell=2.
 \end{cases}
\]
Por tanto cada longitud ponderada por $\log\ell$ es a lo sumo
$\log(2DrX)$. Contamos a lo sumo $r$ torcidos, $X$ primos
y $D$ lugares para cada uno. Esto prueba la cota superior
de \eqref{eq:fu-correccion-global}.

Finalmente, \eqref{eq:fu-k-comparacion} da
$\mathcal K(E,t)=\sum_{\ell,e}W_\ell(E,e)\log\ell$.
La Proposición~\ref{prop:fu-finitud} ya permite sumar las
longitudes de Sha. Sumar \eqref{eq:fu-defecto-local}
prueba la identidad. No se resta una serie infinita de otra.
\end{proof}

La cota es deliberadamente gruesa y no depende del discriminante.
Para grado y corte acotados, controlar uniformemente
$\mathcal V(E,t)$ equivale a controlar
$\mathcal K(E,t)$, con el error explícito anterior.
La fórmula \eqref{eq:fu-correccion-primo} desarrolla y precisa
la corrección de A93, (9.20), p.~73: el coste del torcido aparece
como $v_\ell(e)$. El tamaño de los $K$-grupos sigue siendo
la parte aritmética principal.

\subsection{Un ejemplo donde la corrección no desaparece}

Tomemos $E=\Q(i)$, $\ell=5$ y $e=4$. Como
$5=(2+i)(2-i)$, hay dos lugares $v_1,v_2$ sobre $5$ y
$E_{v_1}=E_{v_2}=\Q_5$. En ambos lugares la imagen ciclotómica
es $\Z_5^\times$. Su parte de orden cuatro es anulada por la
potencia $e=4$ y el nivel principal es uno. Así,
\[
 a_{5,4}(E_{v_1})=a_{5,4}(E_{v_2})=1.
\]
La imagen ciclotómica global también es $\Z_5^\times$:
la intersección de $\Q(i)$ con $\Q(\mu_{5^\infty})$ es $\Q$,
pues la primera extensión cuadrática ramifica en $2$ y la segunda
solo puede ramificar en $5$. Por tanto $a_{5,4}(E)=1$.
La fórmula de dualidad da la igualdad concreta
\begin{equation}
 W_5(\Q(i),4)=s_{5,4}(\Q(i))+1.
 \label{eq:fu-ejemplo-gauss}
\end{equation}
Los grupos son finitos por la comparación clásica. En particular,
$5$ divide $\#K_8(\Z[i])$. El cálculo no determina la longitud
de Sha: demuestra que identificarla directamente con la de $H^2$
perdería una unidad, aun en un cuerpo cuadrático.

\subsection{La longitud exacta a partir de una presentación}

Las cotas por determinantes de C levantan un endomorfismo a un
módulo libre. Las relaciones del módulo permiten recuperar la
longitud exacta, incluso si el determinante de un levantamiento
particular se anula.

\begin{proposicion}[Menores maximales]
\label{prop:fu-menores}
Sea $X$ un $\Z_\ell$-módulo finitamente generado con presentación
\[
 \Z_\ell^a\xrightarrow{P}\Z_\ell^g\xrightarrow{\pi}X
 \longrightarrow0,
\]
y sea $f\in\operatorname{End}_{\Z_\ell}(X)$.
Elijamos una matriz $F\in M_g(\Z_\ell)$ con $\pi F=f\pi$
y formemos la matriz de bloques $A=[\,P\mid F\,]$.
El conúcleo de $f$ es finito si y solo si $A$ tiene rango $g$
sobre $\Q_\ell$. En ese caso,
\begin{equation}
 \length_{\Z_\ell}(X/fX)
 =\min_I v_\ell(\det A_I),
 \label{eq:fu-menores}
\end{equation}
donde $I$ recorre las elecciones de $g$ columnas y
$v_\ell(0)=\infty$. Para $g=0$ se usa el menor vacío igual a uno.
\end{proposicion}

\begin{proof}
La libertad del módulo de presentación permite levantar las
imágenes de sus $g$ vectores básicos, así que existe $F$.
Una preimagen de $fX$ en $\Z_\ell^g$ es
$P\Z_\ell^a+F\Z_\ell^g$. En consecuencia,
\[
 X/fX\simeq\Z_\ell^g/A\Z_\ell^{a+g}.
\]
La forma normal de Smith sobre el anillo de valoración discreta
reduce $A$, mediante operaciones invertibles de filas y columnas,
a una diagonal rectangular con entradas
$\ell^{b_1},\ldots,\ell^{b_q}$, donde $q$ es su rango.
Si $q<g$, el conúcleo tiene parte libre no nula y longitud infinita.
Si $q=g$, su longitud es $b_1+\cdots+b_g$.

El ideal generado por los menores maximales no cambia bajo esas
operaciones: cada menor transformado es una combinación integral
de los anteriores, y las operaciones inversas dan la otra inclusión.
En la diagonal ese ideal es
$(\ell^{b_1+\cdots+b_g})$. En un anillo de valoración discreta,
la valoración de un ideal generado por una colección finita
es el mínimo de las valoraciones de sus generadores.
Esto prueba \eqref{eq:fu-menores}.
\end{proof}

El menor formado solo por $F$ da
$\length(X/fX)\le v_\ell(\det F)$ cuando este último es no nulo.
El mínimo de todos los menores conserva también las relaciones.
Por ejemplo, si $X=\Z_\ell/\ell^a\Z_\ell$ y $f$ es
multiplicación por $\ell^b$, con $a\ge1$, $b\ge0$, la matriz
$[\ell^a\mid\ell^b]$ da longitud $\min(a,b)$.
Si $f=0$, puede tomarse $F=0$ y el mismo procedimiento da
longitud $a$, aunque $\det F=0$.

\subsection{Separar la torsión del cociente libre}

\begin{proposicion}[Dos contribuciones a la longitud]
\label{prop:fu-torsion}
Sea $T=X_{\mathrm{tors}}$ y $F_X=X/T$, con $X$ finitamente
generado sobre $\Z_\ell$. Un endomorfismo $f$ induce
$f_T$ y $\overline f$ en esos módulos. Si
$\det(\overline f)\ne0$, entonces
\begin{equation}
 \length(X/fX)
 =\length(T/f_TT)+v_\ell(\det\overline f).
 \label{eq:fu-torsion-libre}
\end{equation}
Si $F_X=0$, el determinante vacío vale uno.
\end{proposicion}

\begin{proof}
El módulo $F_X$ es libre. El determinante no nulo implica que
$\overline f$ es inyectivo. Si $f(x)\in T$, su imagen en $F_X$
es cero, de modo que $x\in T$. Por tanto
$fX\cap T=f_TT$ y hay una sucesión exacta
\[
 0\longrightarrow T/f_TT\longrightarrow X/fX
 \longrightarrow F_X/\overline f F_X\longrightarrow0.
\]
La forma normal de Smith calcula la longitud del último término
como $v_\ell(\det\overline f)$. La aditividad prueba la fórmula.
\end{proof}

No se exige que una descomposición $X\simeq F_X\oplus T$
sea compatible con $f$. Las contribuciones se separan mediante
la sucesión exacta. Si el determinante en $F_X$ se anula,
$F_X/\overline f F_X$ tiene rango positivo y $X/fX$ no es finito.
Así se distingue una singularidad del operador libre de la
anulación accidental del determinante de un levantamiento redundante.

\subsection{Cuándo una familia de presentaciones tiene cota común}

\begin{proposicion}[Compacidad sin especializaciones singulares]
\label{prop:fu-compacidad}
Fijemos $\ell$, $g$, $b$ y un espacio compacto $S$.
Sea $A:S\to M_{g,b}(\Z_\ell)$ una aplicación continua.
Si $A(s)$ tiene rango $g$ sobre $\Q_\ell$ para todo $s\in S$,
las longitudes de sus conúcleos admiten una cota común.
\end{proposicion}

\begin{proof}
Si $g=0$ todas las longitudes son cero; si $S$ es vacío,
no hay nada que probar. En los demás casos ponemos
\[
 d(s)=\max_I|\det A_I(s)|_\ell.
\]
Es una función continua, positiva, porque hay finitos menores
y en cada punto alguno es no nulo. Su imagen compacta tiene
mínimo $\delta>0$. La Proposición~\ref{prop:fu-menores} da
$\length\operatorname{coker}A(s)=-\log_\ell d(s)
\le-\log_\ell\delta$.
\end{proof}

La exclusión de puntos singulares es decisiva. Con $X=\Z_\ell$
y $f_n$ multiplicación por $\ell^n$, todos los conúcleos son
finitos y de un solo generador, pero sus longitudes son $n$.
Las matrices convergen a la matriz cero, cuyo conúcleo es infinito.
En términos de acciones puede tomarse $\gamma_n=1+\ell^n$ y
$f_n=\gamma_n-1$. La generación acotada, la dimensión fija y
la finitud en cada punto no proporcionan por sí solas la cota común.
El criterio de compacidad requiere además una familia continua
con presentación común y rango completo en todo el compacto.

\subsection{Imágenes exactas y crecimiento por capas}

Pasamos al mecanismo algebraico de C, §5, pp.~7--9.
Escribimos $w\equiv1\pmod{\ell^2}$ cuando
$(w-1)X\subseteq\ell^2X$. Esta es una condición sobre imágenes.
Con torsión no permite dividir automáticamente $w-1$ por
$\ell^2$ dentro de $\operatorname{End}(X)$.

\begin{lema}[Una igualdad de imágenes]
\label{lem:fu-imagen}
Sea $X$ finitamente generado sobre $\Z_\ell$ y sea
$w\in\operatorname{End}(X)$ con $(w-1)X\subseteq\ell^2X$.
Si $Y=(w-1)X$, entonces, para todo $m\ge0$,
\begin{equation}
 (w^{\ell^m}-1)X=\ell^mY.
 \label{eq:fu-imagen-exacta}
\end{equation}
\end{lema}

\begin{proof}
Pongamos $N=w-1$. Como $NX\subseteq\ell^2X$,
$N(Y)=N^2X\subseteq N(\ell^2X)=\ell^2Y$.
Por tanto cada $w^i-1$ lleva $Y$ a $\ell^2Y$.
Para $Q_w=1+w+\cdots+w^{\ell-1}$ se sigue
\[
 (Q_w-\ell)Y\subseteq\ell^2Y,\qquad
 Q_wY\subseteq\ell Y,\qquad
 \ell Y\subseteq Q_wY+\ell^2Y.
\]
El módulo $M=\ell Y/Q_wY$ es finitamente generado y satisface
$M=\ell M$. El lema de Nakayama da $M=0$, de donde
$Q_wY=\ell Y$. La factorización
$w^\ell-1=Q_w(w-1)$ prueba el caso $m=1$.

La imagen obtenida sigue contenida en $\ell^2X$.
Podemos aplicar el mismo paso a $w^\ell$, después a
$w^{\ell^2}$ y así sucesivamente. Partiendo de $m=0$
se obtiene \eqref{eq:fu-imagen-exacta} por inducción.
\end{proof}

\begin{proposicion}[Crecimiento exacto con torsión]
\label{prop:fu-crecimiento}
En las hipótesis del lema, supongamos que $X/Y$ es finito y
escribamos
\[
 Y\simeq\Z_\ell^\rho\oplus
       \bigoplus_{i=1}^{h}\Z_\ell/\ell^{a_i}\Z_\ell,
 \qquad a_i\ge1.
\]
Entonces, para todo $m\ge0$,
\begin{equation}
 \length\operatorname{coker}(w^{\ell^m}-1)
 =\length(X/Y)+m\rho+\sum_{i=1}^{h}\min(m,a_i).
 \label{eq:fu-crecimiento-exacto}
\end{equation}
El aumento al pasar de $m$ a $m+1$ es
$\rho+\#\{i:a_i>m\}$. Para $m\ge\max_i a_i$, con máximo
vacío igual a cero, el crecimiento es lineal de pendiente
$\rho=\operatorname{rank}_{\Z_\ell}X$.
\end{proposicion}

\begin{proof}
Por el lema, el conúcleo es $X/\ell^mY$. La sucesión
\[
 0\longrightarrow Y/\ell^mY\longrightarrow X/\ell^mY
 \longrightarrow X/Y\longrightarrow0
\]
es exacta. El primer término tiene longitud
$m\rho+\sum_i\min(m,a_i)$ por la descomposición de $Y$.
Esto prueba la fórmula y, restando dos niveles sucesivos,
el incremento indicado. Como $X/Y$ es finito, $X$ e $Y$
tienen el mismo rango.
\end{proof}

Si $X$ admite $R$ generadores, todo submódulo suyo admite
a lo sumo $R$: la preimagen del submódulo en $\Z_\ell^R$
es libre de rango a lo sumo $R$ y lo sobreyecta.
En particular $\rho+h\le R$, y la fórmula exacta recupera
el resto $Rm$ de C. También muestra qué controla la etapa
anterior al crecimiento lineal: los exponentes de torsión de $Y$.

Para la acción torcida $g=u^e\gamma$ de C, tomamos
$w=g^{\ell^{j_0}}$, donde
$j_0=\lfloor\log_\ell R\rfloor+3$ si $R>0$, y $j_0=3$
si $R=0$. La unipotencia módulo $\ell$ implica
$(w-1)X\subseteq\ell^2X$: en $X/\ell X$ basta elevar
$g-1$ a una potencia $\ell^k\ge R$, y una elevación adicional
a $\ell$ da la congruencia módulo $\ell^2$.
La fórmula \eqref{eq:fu-crecimiento-exacto} se aplica a
$\nu=j_0+m$ cuando el conúcleo en $j_0$ es finito.
Para $\nu<j_0$, la factorización de potencias da un cociente
del nivel $j_0$. Se conserva así la misma capa de referencia
de la Nota~\ref{sec:cohomologia}.

\subsection{Sumar las distancias mediante un factorial}

La separación ultramétrica de la Nota~\ref{sec:cohomologia}
acotaba cada distancia por $\lfloor\log_\ell Z\rfloor$.
Cuando se necesita la suma de longitudes, conviene sumar
las valoraciones antes de reemplazarlas por un máximo.

\begin{lema}[Suma de distancias a un centro]
\label{lem:fu-factorial}
Sea $c\in\Z_\ell$ y sea $Z\ge1$ entero.
Elijamos $e_c\in\{1,\ldots,Z\}$ que maximice $v_\ell(c-e_c)$,
permitiendo el valor infinito. Entonces
\begin{equation}
 \begin{split}
 \sum_{\substack{1\le e\le Z\\e\ne e_c}}v_\ell(c-e)
 &\le v_\ell\bigl((e_c-1)!(Z-e_c)!\bigr)\\
 &\le v_\ell((Z-1)!)
 =\sum_{k\ge1}\left\lfloor\frac{Z-1}{\ell^k}\right\rfloor
 \le\frac{Z-1}{\ell-1}.
 \end{split}
 \label{eq:fu-factorial}
\end{equation}
\end{lema}

\begin{proof}
Para $e\ne e_c$, la desigualdad ultramétrica y la elección
del máximo dan
\[
 v_\ell(e-e_c)\ge
 \min\{v_\ell(c-e),v_\ell(c-e_c)\}=v_\ell(c-e).
\]
Sumando en todos esos enteros, el producto de las diferencias
en valor absoluto ordinario es $(e_c-1)!(Z-e_c)!$.
Este producto divide $(Z-1)!$, porque el cociente es
$\binom{Z-1}{e_c-1}$.
Para la igualdad siguiente contamos una vez cada múltiplo
de $\ell$, otra vez cada múltiplo de $\ell^2$, y así sucesivamente.
Finalmente omitimos los suelos y sumamos la serie geométrica.
La prueba incluye $Z=1$, con suma vacía y $0!=1$.
\end{proof}

La exclusión de un entero evita la posible valoración infinita.
Si el centro es un entero del intervalo, se excluye precisamente
ese entero. No se supone que los demás enteros estén a una
distancia constante del centro.

\subsection{Un presupuesto para todos los torcidos ordinarios}

Usamos las hipótesis de C: $X$ está generado por $R$ elementos,
$\gamma$ es unipotente módulo $\ell$, y $u\in\Z_\ell^\times$
tiene nivel $a_u=v_\ell(u-1)\ge1$, con $a_u\ge2$ si $\ell=2$.
Escribimos
\[
 \tau(\gamma,u;e,\nu)
 =\length\operatorname{coker}\bigl((u^e\gamma)^{\ell^\nu}-1\bigr).
\]
Los Teoremas A y B de C dan centros $c_1,\ldots,c_R\in\Z_\ell$,
comunes a todos los torcidos y capas de este par $(\gamma,u)$,
y las desigualdades de la Nota~\ref{sec:cohomologia}, §8.17.
Estas cotas se reciben con las hipótesis indicadas.
Para la suma utilizamos la forma detallada del Teorema A,
\[
 \tau(\gamma,u;e,\nu)
 \le R\bigl(\max\{\nu,j_0\}+a_u\bigr)
       +\sum_{i=1}^{R}v_\ell(c_i-e).
\]
Conservar la suma de valoraciones es lo que permite el factor
$R$ del término factorial siguiente.

\begin{proposicion}[Suma ordinaria y término excepcional]
\label{prop:fu-suma-torre}
Para $Z\ge1$, elijamos $e_i\in\{1,\ldots,Z\}$ que maximice
$v_\ell(c_i-e_i)$ para cada centro y pongamos
$\mathcal E_Z=\{e_1,\ldots,e_R\}$.
Con $M_\nu=\max\{\nu,j_0\}+a_u$, se tiene
\begin{equation}
 \sum_{\substack{1\le e\le Z\\e\notin\mathcal E_Z}}
      \tau(\gamma,u;e,\nu)
 \le RZ M_\nu+R\,v_\ell((Z-1)!).
 \label{eq:fu-suma-ordinaria}
\end{equation}
Si los conúcleos en la capa $j_0$ de los torcidos excepcionales
son finitos, todos los términos de la suma completa son finitos y
\begin{equation}
 \begin{split}
 \sum_{e=1}^{Z}\tau(\gamma,u;e,\nu)
 \le{}&RZ M_\nu+R\,v_\ell((Z-1)!)\\
 &+\sum_{e\in\mathcal E_Z}\tau(\gamma,u;e,j_0)
   +R|\mathcal E_Z|(\nu-j_0)_+.
 \end{split}
 \label{eq:fu-suma-completa}
\end{equation}
Para $R=0$ las dos sumas son cero y $\mathcal E_Z$ es vacío.
\end{proposicion}

\begin{proof}
Fuera de $\mathcal E_Z$, la cota por centros de C es finita.
Al sumarla aparecen a lo sumo $RZ$ copias de $M_\nu$.
Para cada centro, las valoraciones restantes son no negativas
porque $c_i\in\Z_\ell$. Por tanto añadir los enteros
excluidos distintos de $e_i$ solo aumenta la suma, y
el Lema~\ref{lem:fu-factorial} la acota por
$v_\ell((Z-1)!)$. Esto prueba \eqref{eq:fu-suma-ordinaria}.
Para cada $e\in\mathcal E_Z$, la segunda cota de C da
\[
 \tau(\gamma,u;e,\nu)
 \le\tau(\gamma,u;e,j_0)+R(\nu-j_0)_+.
\]
Sumarlas y añadir el presupuesto ordinario prueba la segunda
desigualdad y la finitud afirmada.
\end{proof}

La pérdida de proximidad se acota ahora por
$R(Z-1)/(\ell-1)$, en lugar de
$RZ\lfloor\log_\ell Z\rfloor$.
La contribución $RZM_\nu$ conserva el nivel de la torre.
Cuando $\nu=O(\log x)$, esta mejora de la suma de distancias
no elimina por sí sola el término de tamaño $Z\log x$
procedente de la altura.

\subsection{Aplicación a una torre aritmética}

Sea $L_\infty/L_0$ una torre ciclotómica $\Z_\ell$ con
$L_0$ totalmente imaginario, $\mu_\ell\subset L_0$ para
$\ell$ impar y $\mu_4\subset L_0$ para $\ell=2$.
Sea $X_\infty$ el módulo no ramificado y escindido sobre
$\ell$ descrito en C, §7, pp.~9--11.
Elegimos un generador topológico $\sigma$ de la torre;
$\gamma$ es su acción en $X_\infty$ y
$u=\chi_\ell(\sigma)$ genera la imagen ciclotómica.
Las hipótesis de raíces dan el nivel requerido para $u$.
Recibimos la generación de $X_\infty$ por $R$ elementos
sobre $\Z_\ell$, la unipotencia residual de $\gamma$
de C, Proposición~7.2 e Input~1, y el puente de descenso,
con las mismas condiciones locales,
\begin{equation}
 s_{\ell,e}(L_\nu)=\tau(\gamma,u;e,\nu),\qquad e\ge1.
 \label{eq:fu-puente-torre}
\end{equation}
Este es el lugar donde se identifica el módulo algebraico con
cohomología. Las matrices arbitrarias de los ejemplos anteriores
no proporcionan esa identificación.

La Proposición~\ref{prop:fu-finitud}, aplicada a $L_{j_0}$,
prueba la finitud de todos los términos de esa capa y da
\[
 \tau(\gamma,u;e,j_0)
 \le v_\ell\bigl(\#K_{2e}(\mathcal O_{L_{j_0}})\bigr).
\]
En consecuencia, \eqref{eq:fu-suma-completa} se convierte en
\begin{equation}
 \begin{split}
 \sum_{e=1}^{Z}s_{\ell,e}(L_\nu)\log\ell
 \le{}&(\log\ell)
    \left[RZM_\nu+R\,v_\ell((Z-1)!)
                  +R^2(\nu-j_0)_+\right]\\
 &+\sum_{e\in\mathcal E_Z}
       \log\#K_{2e}(\mathcal O_{L_{j_0}}).
 \end{split}
 \label{eq:fu-presupuesto-torre}
\end{equation}
En el último paso se usó $|\mathcal E_Z|\le R$ y se acotó
la parte $\ell$-primaria de cada cardinalidad por la cardinalidad
completa. Para una torre fija y un corte $Z$ fijo, esto es un
presupuesto finito con crecimiento a lo sumo lineal en $\nu$.
El término inicial vive en $L_{j_0}$, no en la capa variable.

Cuando se quiere un presupuesto uniforme en cuerpos variables,
la expresión de teoría $K$ debe ser estimada. A93,
Proposición~9.4, p.~67, formula precisamente una entrada de
otra fuerza:
\begin{equation}
 (\log\ell)W_\ell(L,e)
 \le\exp\!\left((2e[L:\Q])^{C_K}
                       \log(2|\operatorname{disc}L|)\right).
 \label{eq:fu-lemmak}
\end{equation}
Su hipótesis de raíces de unidad y su exponente absoluto se
conservan. Los pasos homológicos, de torsión integral y de
caracteres de clases que la fuente utiliza para demostrarla
se describen en A93, pp.~67--74. La finitud clásica obtenida
aquí permite separar esa cuestión cualitativa de la estimación
de tamaño; no demuestra \eqref{eq:fu-lemmak}.

\subsection{El alcance uniforme y el siguiente desarrollo}

La nueva información puede localizarse en el presupuesto de la
Nota~\ref{sec:cohomologia} sin cambiar sus cuantificadores.

\begin{center}
\small
\begin{tabularx}{\textwidth}{@{}p{0.29\textwidth}YY@{}}
\toprule
Cantidad & Control obtenido & Dato que todavía debe acotarse\\
\midrule
$\mathcal V(E,t)$, $E,t$ fijos
 & Finitud y soporte finito mediante $K_{2e}(\mathcal O_E)$.
 & Tamaño de esos grupos al variar $E,t$.\\
Corrección de dualidad
 & Cota explícita en grado y corte; soporte $\ell\le Dr+1$.
 & Ningún discriminante interviene en esta corrección.\\
Conúcleo presentado
 & Igualdad por menores maximales; separación de torsión y rango.
 & Presentaciones y valoraciones controladas para la familia.\\
Capas de una torre
 & Crecimiento exacto desde una capa suficientemente profunda.
 & Longitud inicial y generación uniforme del módulo.\\
Torcidos ordinarios
 & Suma factorial y a lo sumo $R$ excepciones.
 & Cota de capa fija para los torcidos excepcionales.\\
\bottomrule
\end{tabularx}
\end{center}

En particular, la generación finita del módulo de cada torre
y la anulación de su invariante $\mu$ no eligen un mismo $R$
para todas las bases.
El orden requerido en la rama diádica sigue siendo
\[
 \exists R\ \forall(q,d)\text{ admisibles}\ \exists L_0,
\]
con la cota fijada antes de la pareja auxiliar.
La selección de la ventana de la Nota~\ref{sec:cohomologia}
usa ese mismo control del número de excepciones.
La mejora factorial tampoco sustituye las estimaciones de
los otros primos y caracteres.

La Nota~\ref{sec:capa-fija} desarrolla el paso cuantitativo de A93,
Proposición~9.4, pp.~67--74: Bockstein, conteo por fibras,
cancelación de la altura y absorción de los costes.
Deduce \eqref{eq:fu-lemmak} desde las entradas de homología,
clases y torsión que conserva explícitas. Permanece además
la comparación theta de la Nota~\ref{sec:construccion-selmer}:
construir su filtración, sus mapas, sus multiplicidades y el
control de los núcleos. La finitud del presupuesto de torcidos
de Tate no identifica los Selmer de otros coeficientes ni
construye esos mapas.

\subsection{Ejercicios resueltos}

\begin{ejercicio}\label{ej:fu-h0}
Calcule $a_{\ell,e}(\Q_\ell)$ para $\ell$ impar y para
$\ell=2$. Aplique el resultado a $e=2$ e identifique todos
los primos con invariantes locales no nulos.
\end{ejercicio}

\paragraph{Solución del Ejercicio~\ref{ej:fu-h0}.}
Para $\ell$ impar la imagen es
$\mu_{\ell-1}\times(1+\ell\Z_\ell)$, de modo que
\[
 a_{\ell,e}(\Q_\ell)=
 \begin{cases}
  0,&\ell-1\nmid e,\\
  1+v_\ell(e),&\ell-1\mid e.
 \end{cases}
\]
En $\Q_2$ la imagen es $\{\pm1\}\times(1+4\Z_2)$.
Si $e$ es impar, $-1$ da valoración uno y todos los elementos
tienen potencia impar congruente a uno módulo dos; la longitud
es uno. Si $e$ es par, el signo desaparece y la longitud es
$2+v_2(e)$. Para $e=2$ solo contribuyen $2$ y $3$, con
longitudes $3$ y $1$. El producto de las cardinalidades
locales es $2^3\cdot3=24$.
Este cálculo local no aplica sin más la dualidad global
totalmente imaginaria al cuerpo real $\Q$.

\begin{ejercicio}\label{ej:fu-menores}
Sea $X=\Z_\ell\oplus\Z_\ell/\ell^5\Z_\ell$ y sea
$f(x,\overline y)=(\ell^2x,\ell^3\overline y)$.
Calcule la longitud del conúcleo por una presentación y
por la separación de la torsión.
\end{ejercicio}

\paragraph{Solución del Ejercicio~\ref{ej:fu-menores}.}
Una matriz de presentación del conúcleo es
\[
 \begin{pmatrix}0&\ell^2&0\\\ell^5&0&\ell^3\end{pmatrix}.
\]
Los menores de tamaño dos tienen valoraciones $7$, $\infty$
y $5$. La longitud es $5$. Por otro lado, el cociente libre
contribuye $2$ y el cociente de la torsión por multiplicación
por $\ell^3$ contribuye $\min(5,3)=3$.
Ambos métodos dan $2+3=5$.

\clearpage
\begin{ejercicio}\label{ej:fu-crecimiento}
Suponga que $w-1$ es multiplicación por $\ell^2$ en
$X=\Z_\ell\oplus\Z_\ell/\ell^4\Z_\ell$.
Calcule la longitud de $\operatorname{coker}(w^{\ell^m}-1)$
para todo $m\ge0$ y señale cuándo se estabiliza su incremento.
\end{ejercicio}

\paragraph{Solución del Ejercicio~\ref{ej:fu-crecimiento}.}
Aquí $Y=\ell^2X\simeq\Z_\ell\oplus\Z_\ell/\ell^2\Z_\ell$
y $\length(X/Y)=4$. La fórmula exacta es
\[
 4+m+\min(m,2).
\]
Los primeros valores son $4,6,8,9,10,\ldots$.
El incremento es dos en los dos primeros pasos y uno
desde $m=2$. El rango libre determina el crecimiento final;
la torsión de la imagen determina la transición.

\begin{ejercicio}\label{ej:fu-factorial}
Tome $\ell=3$, $Z=10$ y el centro $c=4$.
Calcule la suma de $v_3(4-e)$ para $1\le e\le10$, $e\ne4$,
y compárela con las dos cotas factoriales.
\end{ejercicio}

\paragraph{Solución del Ejercicio~\ref{ej:fu-factorial}.}
Se excluye $e_c=4$. Entre las diferencias no nulas, solo
$e=1,7,10$ aportan valoración positiva, igual a uno en
cada caso. La suma es $3$ y coincide con
$v_3(3!6!)=1+2=3$. La cota independiente de la posición es
$v_3(9!)=3+1=4$, y la geométrica da $9/2$.
La estimación término a término daría
$10\lfloor\log_3 10\rfloor=20$.

\begin{ejercicio}\label{ej:fu-uniformidad}
En una familia de cuerpos totalmente imaginarios con grado
a lo sumo $D$ y corte fijo $t\ge2$, suponga conocida una
cota común $B$ para $\mathcal V(E,t)$.
¿Qué cota se obtiene para
$\prod_{1\le e<t}\#K_{2e}(\mathcal O_E)$?
¿Por qué la finitud separada de esos productos no basta
para obtener la cota cohomológica requerida en una ventana?
\end{ejercicio}

\paragraph{Solución del Ejercicio~\ref{ej:fu-uniformidad}.}
Con $r=t-1$ y $X=Dr+1$, la Proposición~\ref{prop:fu-presupuesto-exacto}
da
\[
 \prod_{1\le e<t}\#K_{2e}(\mathcal O_E)
 \le\exp\bigl(B+DrX\log(2DrX)\bigr).
\]
En la dirección contraria, una cota para el logaritmo del
producto acota $\mathcal V(E,t)$ directamente.
Una colección de enteros finitos puede ser no acotada;
además, en la aplicación cambian el grado, el discriminante
y el corte. Hace falta una estimación con dependencias y
cuantificadores compatibles con esa variación, como la
entrada de capa fija junto con los controles uniformes
de torre y la selección de ventana.
